\begin{document}

\def\rmd{\mathrm{d}}%
\def\rme{\mathrm{e}}%
\def\rmi{\mathrm{i}}%

\section{Introduction}

The band structure of periodic lattices is a cornerstone of modern
condensed matter physics. In Hermitian systems, the energy spectrum
under open boundary conditions (OBCs) is consistent with the Bloch
bands in the thermodynamic limit. However, non-Hermitian lattices
can exhibit the non-Hermitian skin effect, where the spectrum and
eigenstates under OBC deviate from the Bloch bands and Bloch wave
functions \cite{leeAnomalousEdgeState2016,yaoEdgeStatesTopological2018}.
The deviation of the OBC bands from the Bloch bands has been observed
in various non-Hermitian classical \cite{helbigGeneralizedBulkBoundary2020,hofmannReciprocalSkinEffect2020,weidemannTopologicalFunnelingLight2020,wangGeneratingArbitraryTopological2021,zhangObservationHigherorderNonHermitian2021,zouObservationHybridHigherorder2021,liuComplexSkinModes2022,guTransientNonHermitianSkin2022,wangExtendedStateLocalized2022,shangExperimentalIdentificationSecondOrder2022,liuExperimentalObservationNonHermitian2023,wuEvidencingNonBlochDynamics2023,zhangElectricalCircuitRealization2023,zhuHigherRankChirality2023,liuLocalizationChiralEdge2024,liObservationGaugeField2025}
and quantum \cite{xiaoNonHermitianBulkBoundary2020,wangDetectingNonBlochTopological2021,xiaoObservationNonBlochParityTime2021,zhangObservationNonHermitianTopology2021,linObservationNonHermitianTopological2022,linTopologicalPhaseTransitions2022,chuBroadColossalEdge2023,ochkanNonHermitianTopologyMultiterminal2024,zhaoTwodimensionalNonHermitianSkin2025}
systems.

For one-dimensional (1D) non-Hermitian systems, the thermodynamic
limit of OBC spectra and eigenstates is described by the non-Bloch
band theory using the generalized Brillouin zone (GBZ) \cite{yaoEdgeStatesTopological2018,yokomizoNonBlochBandTheory2019,yangNonHermitianBulkBoundaryCorrespondence2020},
which has been verified in numerous studies \cite{okumaTopologicalPhaseTransition2019,zhuPhotonicNonHermitianSkin2020,xiaoNonHermitianBulkBoundary2020,bartlettIlluminatingBulkboundaryCorrespondence2021,longhiNonHermitianTopologicalPhase2021,wangDetectingNonBlochTopological2021,xiaoObservationNonBlochParityTime2021}.
However, the understanding of non-Hermitian bands in higher dimensions
remains limited due to the geometry dependence of the energy bands
 \cite{zhangUniversalNonHermitianSkin2022, chengTruncationdependent$mathcalPT$Phase2022,fangGeometrydependentSkinEffects2022,wanObservationGeometrydependentSkin2023,zhouObservationGeometrydependentSkin2023,qinGeometrydependentSkinEffect2024}. 
As schematically illustrated
in Fig. \ref{fig:illustration-GD-bands}, for two- and higher-dimensional
non-Hermitian lattices, different geometries ($G_{1}$ and $G_{2}$)
can yield distinct spectra ($\sigma_{1}$ and $\sigma_{2}$) and eigenstates
($\psi_{1}$ and $\psi_{2}$), even in the thermodynamic limit. 
 

While the recent frameworks, including the amoeba formulation \cite{wangAmoebaFormulationNonBloch2024}, the dynamical degeneracy splitting theory \cite{zhangDynamicalDegeneracySplitting2023}, the generalized Fermi surface framework \cite{zhangAlgebraicNonHermitianSkin2025} and those for special cases \cite{yaoNonHermitianChernBands2018,yokomizoNonBlochBandsTwodimensional2023,jiangDimensionalTransmutationNonHermiticity2023, zhangEdgeTheoryNonHermitian2024,huTopologicalOriginNonHermitian2024} provide valuable insights, a unified theory that quantitatively captures the geometry-dependent bands is still lacking.


\begin{figure}
\centering

\includegraphics[scale=0.75]{GDSE-20251014}

\caption{\textbf{Illustration of geometry-dependent non-Hermitian bands.} For
the same system, different geometries ($G_{1}$ and $G_{2}$) yield
distinct spectra ($\sigma_{1}$ and $\sigma_{2}$) and eigenstates
($\psi_{1}$ and $\psi_{2}$), even in the thermodynamic limit.}
\label{fig:illustration-GD-bands} 
\end{figure}

To investigate the mechanism of geometry-dependent bands, we develop
a general formulation based on the strip generalized Brillouin zone
(SGBZ), which describes the energy bands when a non-Hermitian lattice
is extended sequentially along its linearly independent axes to infinity.
We demonstrate that the SGBZ of a non-Hermitian lattice can be dependent
of the sequence of axes, and the competition between incompatible
SGBZs results in geometry-dependent bands. Furthermore, through the
transformation of SGBZs, we derive a criterion that a non-Hermitian
system exhibits geometry-dependent bands if and only if its energy
spectrum or imaginary momentum spectrum has a non-zero area.

\section{Results}

\subsection{SGBZ formulation for non-Hermitian bands}

To obtain the SGBZ, we first extend a lattice along a lattice vector
$\mathbf{a}_{1}$ (named the ``major axis''), forming a strip geometry
with finite width, and then take the width to infinity. As illustrated
in Fig. \ref{fig:SGBZ}a, the lattice confined in the strip geometry
can be viewed as a 1D lattice along $\mathbf{a}_{1}$, whose periodic
unit is a slice of sites along the other lattice vector $\mathbf{a}_{2}$
(named the ``minor axis''). Assume that the momentum-space Hamiltonian
of this 1D lattice is $\mathcal{H}_{L_{2}}(\rme^{\rmi k_{1}})$, where
$L_{2}$ is the width of the strip. When the length of the strip is
sufficiently large, the OBC spectrum tends to the GBZ bands defined
by $|\beta_{1}^{(M)}(E)|=|\beta_{1}^{(M+1)}(E)|$ \cite{yaoEdgeStatesTopological2018,yokomizoNonBlochBandTheory2019}.
Here, $\beta_{1}^{(j)}(E)$ is the $j$-th solution of the eigenvalue
equation $\det[E-\mathcal{H}_{L_{2}}(\beta_{1})]=0$, ordered by $|\beta_{1}^{(i)}|\leq|\beta_{1}^{(j)}|,\forall i<j$,
and $-M$ is the lowest degree of $\beta_{1}$ in $\det[E-\mathcal{H}_{L_{2}}(\beta_{1})]$.

However, $M$ diverges as $L_{2}\rightarrow\infty$. 
To address this problem, we use an equivalent form of the GBZ constraint based on the winding number defined by,

\begin{align}
& W_{\text{q1D}}\left(E,\mu_{1},L_{2}\right) \nonumber \\ 
\equiv & \int_{-\pi}^{\pi}\frac{\rmd\theta_{1}}{2\pi\rmi}\frac{\partial\ln\det F_{L_2}(E,\rme^{\mu_1+\rmi\theta_1})}{\partial\theta_{1}},
\end{align}
where $F_{L_2}(E,\beta_1) \equiv E-\mathcal{H}_{L_{2}}(\beta_1)$. 
  By Cauchy argument principle, $W_{\text{q1D}}(E,\mu_{1},L_{2})$ equals $N_{\text{zeros}}-M$, where $N_{\text{zeros}}$ is the number of zeros of $\det[E-\mathcal{H}_{L_{2}}(\beta_{1})]$ with modulii less than $\exp(\mu_1)$. Therefore, as illustrated in Fig.~\ref{fig:SGBZ}b, when $|\beta_1^{(M)}|=|\beta_1^{(M+1)}|$, $W_{\text{q1D}}(E,\mu_1,L_2)$ is negative (light green region) when $\mu_1 < \mu_{1,0}\equiv |\beta_1^{(M)}|$, and positive (light orange region) when $\mu_1 > \mu_{1,0}$. Otherwise, when the GBZ constraint is not satisfied, the regions with positive and negative $W_{\text{q1D}}$ are disconnected.


\begin{figure}[htb]
\centering

\includegraphics[scale=0.75]{SGBZ-formulation-20260325}

\caption{\textbf{Strip geometry and SGBZ constraint.} \textbf{a} Schematic diagram of the
strip geometry of a 2D non-Hermitian lattice, where the cyan region
represents the periodic unit of the strip. 
 \textbf{b} GBZ constraint for the strip geometry with finite width, where the blue and purple dots denote the solutions of the eigenvalue equation. The limit distribution of the solutions is represented by the blue dashed curves. \textbf{c} Illustration of the base manifold
$X(E,\mu_{1})$, where the blue solid lines and cyan dashed lines
represent the $M_{2}$-th and $(M_{2}+1)$-th solutions of the eigenvalue
equations, respectively, and the orange solid curves represent the
winding loops for $w_{1}(\theta_{2};E,\mu_{1})$. 
\textbf{d} SGBZ constraint for the strip geometry with infinite width, where the blue solid lines denote the limit distribution of the solutions of the eigenvalue equation.
}
\label{fig:SGBZ} 
\end{figure}


While $W_{\text{q1D}}$ also diverges as $L_2 \rightarrow \infty$, $W_{\text{q1D}}/L_2$ remains finite. 
The limit of $W_{\text{q1D}}/L_2$ can be calculated from momentum-space Hamiltonian 
$h(\beta_1,\beta_2)$.
We consider
the solutions $\beta_{2}^{(j)}(E,\beta_{1}),j=1,2,\dots,M_{2}+N_{2}$
of $\det[E-h(\beta_{1},\beta_{2})]=0$, where $-M_{2}$ and $N_{2}$
are the lowest and highest degrees of $\beta_{2}$ in $\det[E-h(\beta_{1},\beta_{2})]$,
respectively, and the solutions are ordered by $|\beta_{2}^{(j)}|\leq|\beta_{2}^{(k)}|,\forall j<k$.
For a given radius $|\beta_{1}|=\rme^{\mu_{1}}$,
we define the radius function  $\rho_{2,0}(\theta_{1};E,\mu_{1})$ as,
\begin{equation}
\mu_{2}^{(M_{2})}\leq\rho_{2,0}\leq\mu_{2}^{(M_{2}+1)},\label{eq:def-radius-func}
\end{equation}
where $\mu_{2}^{(j)}(\theta_{1};E,\mu_{1})\equiv\ln|\beta_{2}^{(j)}(E,\rme^{\mu_{1}+\rmi\theta_{1}})|$
are the imaginary momentum components of the $j$-th solutions. We
require that $\mu_{2,0}$ is periodic in $\theta_{1}$, and the equality
in Eq. (\ref{eq:def-radius-func}) holds if and only if  $\mu_{2}^{(M_{2})}(\theta_{1};E,\mu_{1})=\rho_{2,0}(\theta_{1};E,\mu_{1})=\mu_{2}^{(M_{2}+1)}(\theta_{1};E,\mu_{1})$.
Using the radius function, we define the base manifold $X(E,\mu_{1})$
as, 
\begin{align}
X(E,\mu_{1})\equiv & \left\{ \left(\beta_{1},\beta_{2}\right)\in\mathbb{C}^{2}\mid\beta_{1}=\rme^{\mu_{1}+\rmi\theta_{1}},\right.\nonumber \\
 & \left.\beta_{2}=\rme^{\rho_{2,0}(\theta_{1};E,\mu_{1})+\rmi\theta_{2}},\theta_{1},\theta_{2}\in\left[-\pi,\pi\right]\right\} .\label{eq:base-manifold}
\end{align}
As shown in Fig. \ref{fig:SGBZ}c, in the three-dimensional space
defined by $\ln|\beta_{1}|=\mu_{1}$, the solutions $\beta_{2}^{(M_{2}+1)}$
(blue solid curves) lie outside $X(E,\mu_{1})$, while $\beta_{2}^{(M_{2})}$
(cyan dashed curves) are enclosed within $X(E,\mu_{1})$.  With the above preparation, $W_{\text{q1D}}/L_2$ converges to the average major-axis winding number on the base manifold (see Supplementary Information Sec.~S1 for the proof), i.e., 
\begin{equation}
W\left(E,\mu_{1}\right)\equiv\int_{-\pi}^{\pi}\frac{\rmd\theta_{2}}{2\pi}w_{1}\left(\theta_{2};E,\mu_{1}\right),
\end{equation}
where $w_{1}(\theta_{2};E,\mu_{1})$ is the winding number of $\det[E-h(\beta_{1},\beta_{2})]$
when $(\beta_{1},\beta_{2})$ traverses the loop defined by $\text{Arg}(\beta_{2})=\theta_{2}$
on the base manifold $X(E,\mu_{1})$, shown as the orange curves in
Fig. \ref{fig:SGBZ}c.  
$$
\text{}
$$


Therefore, replacing $W_{\text{q1D}}$ with $W$ yields a well-defined GBZ constraint for $L_2 \rightarrow \infty$. 
As shown in Fig.~\ref{fig:SGBZ}d, for each reference energy $E$, if the regions with $W(E,\mu_1)>0$ and $W(E,\mu_1)<0$ are connected at the boundary $\mu_1 = \mu_{1,0}$, then the zeros of $\det[E-h(\beta_{1},\beta_{2})]$ must intersect $X(E,\mu_{1,0})$ (see Supplementary Information Sec.~S1.3 for proof).
By traversing $E$ over the complex plane, these zeros form a 2D subspace of $(\beta_1,\beta_2)\in \mathbb{C}^2$, which we term the SGBZ of a 2D system.
Similar to the 1D GBZs, on the SGBZ, $|\beta_1|$ and $|\beta_2|$ depict the skin localization of the eigenstates, and the energy bands depict the OBC energy spectra.

The SGBZ formulations can also be extended
to higher dimensions. For an $n$-dimensional lattice, a strip is
specified by a sequence of axes $(\mathbf{a}_{1},\mathbf{a}_{2},\dots,\mathbf{a}_{n})$,
and the  average major-axis winding numbers are defined recursively. A detailed
discussion is available in Supplementary Information Sec.~S1.4.



In the SGBZ formulation presented above, it is noted that the status
of $\beta_{1}$ and $\beta_{2}$ is not symmetric. This observation
implies that the SGBZs calculated under different major and minor
axes can be incompatible. In fact, the SGBZ is independent of the
minor axes (see Supplementary Information Sec. S2), but depends on
the selection of the major axes. As an example, we consider the 2D
Hatano-Nelson (HN) model shown in Fig. \ref{fig:OBC}a. The momentum-space
Hamiltonian of the model reads, 
\begin{equation}
h\left(\beta_{x},\beta_{y}\right)=J_{x1}\beta_{x}^{-1}+J_{x2}\beta_{x}+J_{y1}\beta_{y}^{-1}+J_{y2}\beta_{y}.\label{eq:HN}
\end{equation}
As illustrated in Fig. \ref{fig:OBC}b, we consider three different
strips: the {[}10{]}-strip, {[}01{]}-strip, and {[}11{]}-strip, defined
by the major axes $\mathbf{a}_{x}=(1,0)$, $\mathbf{a}_{y}=(0,1)$,
and $\mathbf{a}_{[11]}=(1,1)$, respectively. In Supplementary Information
Sec. S3, we compute the SGBZs for these three distinct strips. The
SGBZs of the {[}10{]}-strip and {[}01{]}-strip are identical, which
reads, 
\begin{align}
\beta_{x} & =\exp\left(\gamma_{x}+\rmi\theta_{x}\right),\quad\beta_{y}=\exp\left(\gamma_{y}+\rmi\theta_{y}\right),\label{eq:xy-SGBZ-HN}
\end{align}
where $\gamma_{\alpha}\equiv\ln(|J_{\alpha1}/J_{\alpha2}|)/2$, $\alpha=x,y$.
The only distinction between the {[}10{]}-SGBZ and {[}01{]}-SGBZ lies
in whether $\mathbf{a}_{x}$ or $\mathbf{a}_{y}$ is the major axis.
For comparison, the SGBZ for the {[}11{]}-strip is given by, 
\begin{equation}
\begin{cases}
\tilde{\beta}_{[11]}=\rme^{\gamma_{x}+\gamma_{y}+\rmi\theta_{[11]}},\\
\tilde{\beta}_{y}=\rme^{\gamma_{y}+\rmi\theta_{y}}\sqrt{\left|\frac{J_{x}^{*}\rme^{\rmi\Delta_{xy}+\rmi\theta_{[11]}}+J_{y}}{J_{x}\rme^{\rmi\Delta_{xy}-\rmi\theta_{[11]}}+J_{y}^{*}}\right|},
\end{cases}\label{eq:11-SGBZ-HN}
\end{equation}
where $\Delta_{xy}\equiv\delta_{x}-\delta_{y}$, $\delta_{\alpha}\equiv\text{Arg}(J_{\alpha1}J_{\alpha2})/2$,
and $J_{\alpha}\equiv J_{\alpha1}/\exp(\gamma_{\alpha}+\rmi\delta_{\alpha})$
for $\alpha=x,y$. If the SGBZs of the three strips are compatible,
the coordinate transformation of momenta, which reads, 
\begin{equation}
\tilde{\beta}_{[11]}=\beta_{x}\beta_{y},\quad\tilde{\beta}_{y}=\beta_{y},\label{eq:trans}
\end{equation}
should hold for Eqs. (\ref{eq:xy-SGBZ-HN}) and (\ref{eq:11-SGBZ-HN}).
However, Eq. (\ref{eq:trans}) fails when $\sin\Delta_{xy}\neq0$,
indicating that the {[}11{]}-SGBZ differs from the {[}10{]}-SGBZ or
{[}01{]}-SGBZ in general. In fact, as shown in Fig. \ref{fig:OBC}c
and \ref{fig:OBC}d, the spectrum of the {[}10{]}-SGBZ ($\sigma_{[10]}$)
or {[}01{]}-SGBZ ($\sigma_{[01]}$) deviates from that of the {[}11{]}-SGBZ
($\sigma_{[11]}$), except when $\rme^{\rmi\delta_{x}}$ is collinear
with $\rme^{\rmi\delta_{y}}$, i.e., $\sin\Delta_{xy}=0$.


\begin{figure*}
\centering

\includegraphics[scale=0.75]{OBC-comparison-20250821}

\caption{\textbf{Relation between SGBZs and geometry-dependent bands.} \textbf{a}
Illustration of the 2D HN model. \textbf{b} Illustration of the {[}10{]}-strip,
{[}01{]}-strip, and {[}11{]}-strip, where the major axes are $\mathbf{a}_{x}=(1,0)$,
$\mathbf{a}_{y}=(0,1)$, and $\mathbf{a}_{[11]}=(1,1)$, respectively.
\textbf{c}, \textbf{d} Spectra of (c) $[10]$-SGBZ or $[01]$-SGBZ,
and (d) $[11]$-SGBZ. \textbf{e}, \textbf{f} Comparison between the
SGBZ bands and the finite-size OBC spectra in parallelogram regions
with different aspect ratios, where (e) illustrates the region with
compatible SGBZs, and (f) illustrates the region with incompatible
SGBZs. The coupling coefficients are $J_{x1}=1+\protect\rmi$, $J_{x2}=1.5+1.2\protect\rmi$,
$J_{y1}=-1+\protect\rmi$, and $J_{y2}=-1.2-0.5\protect\rmi$. In
numerical calculations, the total number of sites is set to be $12800$
(or the nearest integer to $12800$).}
\label{fig:OBC} 
\end{figure*}


To investigate the relation between the compatibility of SGBZs and
the geometry-dependent bands, we numerically calculate the OBC spectrum
in finite-size parallelogram regions with different aspect ratios.
For the compatible case, as shown in Fig. \ref{fig:OBC}e, the parallelogram
region is spanned by the $[10]$-axis and $[01]$-axis, along which
the SGBZs are identical. For all aspect ratios, the OBC spectra match
well with the SGBZ bands. For the incompatible case, as shown in Fig.
\ref{fig:OBC}f, the parallelogram region is spanned by the $[11]$-axis
and $[01]$-axis. In contrast to the compatible case, the OBC spectra
vary with the aspect ratios, and tend to the corresponding SGBZ bands
when the length of one side is much larger than the other side (see Supplementary Information Sec.~S4 for more numerical results). 

To understand this phenomenon, we return to the definition of the
SGBZ. Since the SGBZ is constructed by taking sequential limits, the
strip length remains much larger than the strip width during the process
of taking the width to infinity. Consequently, the SGBZ bands correspond
to the OBC spectra in the limit of both infinite sizes and extreme
aspect ratios. Therefore, when a lattice holds incompatible SGBZs,
the effects of lattice extensions along different axes compete with
each other, preventing the spectrum from converging as the system
size increases (see Supplementary Information Sec.~S5). Unlike Hermitian
systems, where the OBC spectra with different boundary geometries
converge to the Bloch bands, in non-Hermitian systems with incompatible
SGBZs, there is no uniform thermodynamic limit of OBC spectra, resulting
in the geometry-dependent bands.


It is noted that our framework is not limited to the single-band model. For models with multiple bands (see Supplementary Information Sec.~S6) and long-range hoppings (see Supplementary Information Sec.~S7), if $h(\beta_1,\beta_2)$ is holomorphic in the neighborhood of the SGBZ, our framework still holds.


\subsection{Criterion for geometry-dependent bands}

Because the geometry-dependent bands result from the incompatible
SGBZs, a universal criterion for geometry-dependent bands can be derived
by checking the compatibility of SGBZs. Assume that a non-Hermitian
system exhibits uniform bands, i.e., the OBC spectra of the system
converge to a uniform thermodynamic limit. Under the coordinate transformation,
\begin{equation}
\begin{pmatrix}\tilde{\mathbf{a}}_{1} & \tilde{\mathbf{a}}_{2}\end{pmatrix}=\begin{pmatrix}\mathbf{a}_{1} & \mathbf{a}_{2}\end{pmatrix}\mathbf{P},
\end{equation}
the complex momenta of the SGBZs corresponding to $(\tilde{\mathbf{a}}_{1},\tilde{\mathbf{a}}_{2})$
and $(\mathbf{a}_{1},\mathbf{a}_{2})$ must satisfy, 
\begin{equation}
\begin{pmatrix}\ln\tilde{\beta}_{1} & \ln\tilde{\beta}_{2}\end{pmatrix}=\begin{pmatrix}\ln\beta_{1} & \ln\beta_{2}\end{pmatrix}\mathbf{P},\label{eq:trans-beta}
\end{equation}
where $\mathbf{P}\in\mathbb{Z}^{2\times2}$ is an arbitrary transformation
matrix.

Under the basis $(\mathbf{a}_{1},\mathbf{a}_{2})$, as illustrated
in Fig. \ref{fig:uniform-cond}a, we consider the curves $\mu_{2}^{(j)}(\theta_{1};E,\mu_{1,0})$
for $j=M_{2}$ and $M_{2}+1$, where $\mu_{1,0}$ is the critical
point at which $W(E,\mu_{1})$ changes sign. By definition, the common
points of the two curves correspond to SGBZ points with eigenenergy
$E$. First, we consider the permutation of two axes, i.e., $P_{11}=P_{22}=0$
and $P_{12}=P_{21}=1$, where $P_{ij}$ is the matrix element of $\mathbf{P}$.
By the definition of $X(E,\mu_{1,0})$, $\ln|\beta_{1}|=\mu_{1,0}$
is a constant in the original SGBZ. To ensure the transformed points
also satisfy the SGBZ constraint, $\ln|\tilde{\beta}_{1}|=\ln|\beta_{2}|$
should also be constant. Therefore, for uniform bands, the common
points between $\mu_{2}^{(M_{2})}$ and $\mu_{2}^{(M_{2}+1)}$ should
have constant value, which rules out the case i in Fig. \ref{fig:uniform-cond}a.

Second, for a given value of $\beta_{1}$ on the SGBZ, both solutions
$(\beta_{1},\beta_{2}^{(M_{2})})$ and $(\beta_{1},\beta_{2}^{(M_{2}+1)})$
belong to the SGBZ, so the SGBZ points occur in pairs. However, the
transformation in Eq. (\ref{eq:trans-beta}) may disrupt this pairing.
To ensure the transformed points also form pairs, for each SGBZ point
$(\beta_{1},\beta_{2})$, there must exist another SGBZ point $(\beta_{1}^{\prime},\beta_{2}^{\prime})$
at the same eigenenergy such that $\tilde{\beta}_{1}=\tilde{\beta}_{1}^{\prime}$,
that is, 
\begin{equation}
\left(\frac{\beta_{1}}{\beta_{1}^{\prime}}\right)^{P_{11}}=\left(\frac{\beta_{2}^{\prime}}{\beta_{2}}\right)^{P_{21}}.\label{eq:beta-pairing}
\end{equation}
Since the solutions of Eq. (\ref{eq:11-SGBZ-HN}) for $(\beta_{1}^{\prime},\beta_{2}^{\prime})$
depend on $P_{11}/P_{21}$, when $\mathbf{P}$ ranges over all possible
transformation matrices, there must exist infinitely many SGBZ points
paired with $(\beta_{1},\beta_{2})$. Therefore, if the SGBZ points
form a finite set for some eigenenergy, as illustrated in panel ii
of Fig. \ref{fig:uniform-cond}a, the system cannot exhibit uniform
bands.

Third, as illustrated in panel iii of Fig. \ref{fig:uniform-cond}a,
if a horizontal line (black dotted line) passes through both the upper
region of $\mu_{2}^{(M_{2}+1)}$ and the lower region of $\mu_{2}^{(M_{2})}$,

the two roots intersecting with the same horizontal line can be rotated to the same $\tilde{\theta}_1$ via a certain coordinate transformation (see Supplementary Information Sec.~S8). Under such transformation, this case is equivalent to the case of panel i. Thus, the system cannot exhibit uniform bands either.


As discussed above, only when an SGBZ has infinitely many SGBZ points
with constant values of $\mu_{1}$ and $\mu_{2}$ for every eigenenergy
$E$, as illustrated in panel iv of Fig. \ref{fig:uniform-cond}a,
can the system exhibit uniform non-Hermitian bands. In this case,
the base manifold $X(E,\mu_{1,0})$ becomes the (hyper)plane $(\mu_{1},\mu_{2})=(\mu_{1,0},\mu_{2,0})$.
 
To ensure that the sign of $W(E,\mu_1)$ changes at $\mu_{1,0}$, the winding number of the loops in the neighborhood of the plane $\boldsymbol{\mu} = (\mu_{1,0},\mu_{2,0})$ must satisfy the condition shown in Fig.~\ref{fig:uniform-cond}b:
For an arbitrary loop on
the base manifold (orange lines), if the loop intersects the SGBZ
points (blue curve), the winding number is ill-defined. In this case,
as illustrated in the right panel, we can increase or decrease $\mu_{2}$
(or equivalently $\mu_{1}$) to avoid the intersections. For all possible
perturbations, uniform bands require that the maximum and minimum
winding numbers satisfy $w_{\text{max}}=-w_{\text{min}}$. The condition
in Fig. \ref{fig:uniform-cond}b is also sufficient for uniform bands (see Supplementary Information Sec.~S8).

Moreover, the criterion for uniform or geometry-dependent bands also
manifests in the spectrum and the imaginary momentum spectrum. On
one hand, the SGBZ points corresponding to a fixed energy form 1D
curves, so the dimensionality of a 2D uniform spectrum is less than
that of the 2D SGBZ. Therefore, the uniform spectrum should have zero
area. On the other hand, for uniform bands, the SGBZ points corresponding
to a fixed energy should have constant $\mu_{1}$ and $\mu_{2}$,
so the imaginary momentum spectra, i.e., the plots of $\text{Re}(E)$-$\mu_{j}$,
should have zero area. The conclusion can be generalized to higher
dimensions (see Supplementary Information Sec.~S8). For an $n$-dimensional
system, as shown in Fig. \ref{fig:uniform-cond}c, if and only if
both the spectrum and all the imaginary momentum spectra have zero
area, the system exhibits uniform bands. It is noted that zero area
of the SGBZ spectrum is not sufficient for uniform bands; an example
is available in Supplementary Information Sec.~S9. 


\begin{figure}
\centering

\includegraphics[scale=0.75]{illustration-uniform-band-cond-20251014}

\caption{\textbf{Criterion for uniform or geometry-dependent bands.} \textbf{a}
 Minor-axis root distributions of the eigenvalue equation for cases with and without uniform bands, where the blue solid curve and
cyan dashed curve represent $\mu_{2}^{(M_{2}+1)}(\theta_{1};E,\mu_{1,0})$
and $\mu_{2}^{(M_{2})}(\theta_{1};E,\mu_{1,0})$, respectively, and
the black dotted lines in ii and iii are reference lines. In cases
i--iii, uniform bands are not allowed, while in case iv, uniform
bands are allowed. \textbf{b} Requirements on the winding numbers
when a system exhibits a uniform band. The blue curve represents the
solutions of $\det[E-h(\beta_{1},\beta_{2})]=0$ on the plane $(\mu_{1},\mu_{2})=(\mu_{1,0},\mu_{2,0})$,
and the orange curve represents the winding loop. \textbf{c} Illustration
of the criterion for uniform or geometry-dependent bands. For an arbitrary
SGBZ, if the complex energy spectrum and all the imaginary momentum
spectra have zero area, the bands are uniform. Otherwise, if the energy
spectrum or at least one of the imaginary momentum spectra has nonzero
area, the bands are geometry-dependent.}
\label{fig:uniform-cond} 
\end{figure}

% \bibliography{2D-skin-effect}% common bib file
%% if required, the content of .bbl file can be included here once bbl is generated
\begin{thebibliography}{10}
\expandafter\ifx\csname url\endcsname\relax
  \def\url#1{\burl{#1}}\fi
\expandafter\ifx\csname urlprefix\endcsname\relax\def\urlprefix{URL }\fi
\providecommand{\bibinfo}[2]{#2}
\providecommand{\eprint}[2][]{\url{#2}}
\providecommand{\doi}[1]{\url{https://doi.org/#1}}
\bibcommenthead

\bibitem{leeAnomalousEdgeState2016}
\bibinfo{author}{Lee, T.~E.}
\newblock \bibinfo{title}{Anomalous {{Edge State}} in a {{Non-Hermitian Lattice}}}.
\newblock \emph{\bibinfo{journal}{Phys. Rev. Lett.}} \textbf{\bibinfo{volume}{116}}, \bibinfo{pages}{133903} (\bibinfo{year}{2016}).

\bibitem{yaoEdgeStatesTopological2018}
\bibinfo{author}{Yao, S.} \& \bibinfo{author}{Wang, Z.}
\newblock \bibinfo{title}{Edge {{States}} and {{Topological Invariants}} of {{Non-Hermitian Systems}}}.
\newblock \emph{\bibinfo{journal}{Phys. Rev. Lett.}} \textbf{\bibinfo{volume}{121}}, \bibinfo{pages}{086803} (\bibinfo{year}{2018}).

\bibitem{helbigGeneralizedBulkBoundary2020}
\bibinfo{author}{Helbig, T.} \emph{et~al.}
\newblock \bibinfo{title}{Generalized bulk--boundary correspondence in non-{{Hermitian}} topolectrical circuits}.
\newblock \emph{\bibinfo{journal}{Nat. Phys.}} \textbf{\bibinfo{volume}{16}}, \bibinfo{pages}{747--750} (\bibinfo{year}{2020}).

\bibitem{hofmannReciprocalSkinEffect2020}
\bibinfo{author}{Hofmann, T.} \emph{et~al.}
\newblock \bibinfo{title}{Reciprocal skin effect and its realization in a topolectrical circuit}.
\newblock \emph{\bibinfo{journal}{Phys. Rev. Res.}} \textbf{\bibinfo{volume}{2}}, \bibinfo{pages}{023265} (\bibinfo{year}{2020}).

\bibitem{weidemannTopologicalFunnelingLight2020}
\bibinfo{author}{Weidemann, S.} \emph{et~al.}
\newblock \bibinfo{title}{Topological funneling of light}.
\newblock \emph{\bibinfo{journal}{Science}} \textbf{\bibinfo{volume}{368}}, \bibinfo{pages}{311--+} (\bibinfo{year}{2020}).

\bibitem{wangGeneratingArbitraryTopological2021}
\bibinfo{author}{Wang, K.} \emph{et~al.}
\newblock \bibinfo{title}{Generating arbitrary topological windings of a non-{{Hermitian}} band}.
\newblock \emph{\bibinfo{journal}{Science}} \textbf{\bibinfo{volume}{371}}, \bibinfo{pages}{1240--1245} (\bibinfo{year}{2021}).

\bibitem{zhangObservationHigherorderNonHermitian2021}
\bibinfo{author}{Zhang, X.}, \bibinfo{author}{Tian, Y.}, \bibinfo{author}{Jiang, J.-H.}, \bibinfo{author}{Lu, M.-H.} \& \bibinfo{author}{Chen, Y.-F.}
\newblock \bibinfo{title}{Observation of higher-order non-{{Hermitian}} skin effect}.
\newblock \emph{\bibinfo{journal}{Nat. Commun.}} \textbf{\bibinfo{volume}{12}}, \bibinfo{pages}{5377} (\bibinfo{year}{2021}).

\bibitem{zouObservationHybridHigherorder2021}
\bibinfo{author}{Zou, D.} \emph{et~al.}
\newblock \bibinfo{title}{Observation of hybrid higher-order skin-topological effect in non-{{Hermitian}} topolectrical circuits}.
\newblock \emph{\bibinfo{journal}{Nat. Commun.}} \textbf{\bibinfo{volume}{12}}, \bibinfo{pages}{7201} (\bibinfo{year}{2021}).

\bibitem{liuComplexSkinModes2022}
\bibinfo{author}{Liu, Y. G.~N.} \emph{et~al.}
\newblock \bibinfo{title}{Complex skin modes in non-{{Hermitian}} coupled laser arrays}.
\newblock \emph{\bibinfo{journal}{Light-Sci. Appl.}} \textbf{\bibinfo{volume}{11}}, \bibinfo{pages}{336} (\bibinfo{year}{2022}).

\bibitem{guTransientNonHermitianSkin2022}
\bibinfo{author}{Gu, Z.} \emph{et~al.}
\newblock \bibinfo{title}{Transient non-{{Hermitian}} skin effect}.
\newblock \emph{\bibinfo{journal}{Nat. Commun.}} \textbf{\bibinfo{volume}{13}}, \bibinfo{pages}{7668} (\bibinfo{year}{2022}).

\bibitem{wangExtendedStateLocalized2022}
\bibinfo{author}{Wang, W.}, \bibinfo{author}{Wang, X.} \& \bibinfo{author}{Ma, G.}
\newblock \bibinfo{title}{Extended {{State}} in a {{Localized Continuum}}}.
\newblock \emph{\bibinfo{journal}{Phys. Rev. Lett.}} \textbf{\bibinfo{volume}{129}}, \bibinfo{pages}{264301} (\bibinfo{year}{2022}).

\bibitem{shangExperimentalIdentificationSecondOrder2022}
\bibinfo{author}{Shang, C.} \emph{et~al.}
\newblock \bibinfo{title}{Experimental {{Identification}} of the {{Second-Order Non-Hermitian Skin Effect}} with {{Physics-Graph-Informed Machine Learning}}}.
\newblock \emph{\bibinfo{journal}{Adv. Sci.}} \textbf{\bibinfo{volume}{9}}, \bibinfo{pages}{2202922} (\bibinfo{year}{2022}).

\bibitem{liuExperimentalObservationNonHermitian2023}
\bibinfo{author}{Liu, B.} \emph{et~al.}
\newblock \bibinfo{title}{Experimental observation of non-{{Hermitian}} higher-order skin interface states in topological electric circuits}.
\newblock \emph{\bibinfo{journal}{Phys. Rev. Res.}} \textbf{\bibinfo{volume}{5}}, \bibinfo{pages}{043034} (\bibinfo{year}{2023}).

\bibitem{wuEvidencingNonBlochDynamics2023}
\bibinfo{author}{Wu, M.} \emph{et~al.}
\newblock \bibinfo{title}{Evidencing non-{{Bloch}} dynamics in temporal topolectrical circuits}.
\newblock \emph{\bibinfo{journal}{Phys. Rev. B}} \textbf{\bibinfo{volume}{107}}, \bibinfo{pages}{064307} (\bibinfo{year}{2023}).

\bibitem{zhangElectricalCircuitRealization2023}
\bibinfo{author}{Zhang, H.}, \bibinfo{author}{Chen, T.}, \bibinfo{author}{Li, L.}, \bibinfo{author}{Lee, C.~H.} \& \bibinfo{author}{Zhang, X.}
\newblock \bibinfo{title}{Electrical circuit realization of topological switching for the non-{{Hermitian}} skin effect}.
\newblock \emph{\bibinfo{journal}{Phys. Rev. B}} \textbf{\bibinfo{volume}{107}}, \bibinfo{pages}{085426} (\bibinfo{year}{2023}).

\bibitem{zhuHigherRankChirality2023}
\bibinfo{author}{Zhu, P.}, \bibinfo{author}{Sun, X.-Q.}, \bibinfo{author}{Hughes, T.~L.} \& \bibinfo{author}{Bahl, G.}
\newblock \bibinfo{title}{Higher rank chirality and non-{{Hermitian}} skin effect in a topolectrical circuit}.
\newblock \emph{\bibinfo{journal}{Nat. Commun.}} \textbf{\bibinfo{volume}{14}}, \bibinfo{pages}{720} (\bibinfo{year}{2023}).

\bibitem{liuLocalizationChiralEdge2024}
\bibinfo{author}{Liu, G.-G.} \emph{et~al.}
\newblock \bibinfo{title}{Localization of {{Chiral Edge States}} by the {{Non-Hermitian Skin Effect}}}.
\newblock \emph{\bibinfo{journal}{Phys. Rev. Lett.}} \textbf{\bibinfo{volume}{132}}, \bibinfo{pages}{113802} (\bibinfo{year}{2024}).

\bibitem{liObservationGaugeField2025}
\bibinfo{author}{Li, Y.} \emph{et~al.}
\newblock \bibinfo{title}{Observation of gauge field induced non-{{Hermitian}} helical spin skin effects} (\bibinfo{year}{2025}).
\newblock \bibinfo{eprint}{{\href{https://arxiv.org/abs/2504.18063}{{arXiv:2504.18063}}}}.

\bibitem{xiaoNonHermitianBulkBoundary2020}
\bibinfo{author}{Xiao, L.} \emph{et~al.}
\newblock \bibinfo{title}{Non-{{Hermitian}} bulk--boundary correspondence in quantum dynamics}.
\newblock \emph{\bibinfo{journal}{Nat. Phys.}} \textbf{\bibinfo{volume}{16}}, \bibinfo{pages}{761--766} (\bibinfo{year}{2020}).

\bibitem{wangDetectingNonBlochTopological2021}
\bibinfo{author}{Wang, K.} \emph{et~al.}
\newblock \bibinfo{title}{Detecting {{Non-Bloch Topological Invariants}} in {{Quantum Dynamics}}}.
\newblock \emph{\bibinfo{journal}{Phys. Rev. Lett.}} \textbf{\bibinfo{volume}{127}}, \bibinfo{pages}{270602} (\bibinfo{year}{2021}).

\bibitem{xiaoObservationNonBlochParityTime2021}
\bibinfo{author}{Xiao, L.} \emph{et~al.}
\newblock \bibinfo{title}{Observation of {{Non-Bloch Parity-Time Symmetry}} and {{Exceptional Points}}}.
\newblock \emph{\bibinfo{journal}{Phys. Rev. Lett.}} \textbf{\bibinfo{volume}{126}}, \bibinfo{pages}{230402} (\bibinfo{year}{2021}).

\bibitem{zhangObservationNonHermitianTopology2021}
\bibinfo{author}{Zhang, W.} \emph{et~al.}
\newblock \bibinfo{title}{Observation of {{Non-Hermitian Topology}} with {{Nonunitary Dynamics}} of {{Solid-State Spins}}}.
\newblock \emph{\bibinfo{journal}{Phys. Rev. Lett.}} \textbf{\bibinfo{volume}{127}}, \bibinfo{pages}{090501} (\bibinfo{year}{2021}).

\bibitem{linObservationNonHermitianTopological2022}
\bibinfo{author}{Lin, Q.} \emph{et~al.}
\newblock \bibinfo{title}{Observation of non-{{Hermitian}} topological {{Anderson}} insulator in quantum dynamics}.
\newblock \emph{\bibinfo{journal}{Nat. Commun.}} \textbf{\bibinfo{volume}{13}}, \bibinfo{pages}{3229} (\bibinfo{year}{2022}).

\bibitem{linTopologicalPhaseTransitions2022}
\bibinfo{author}{Lin, Q.} \emph{et~al.}
\newblock \bibinfo{title}{Topological {{Phase Transitions}} and {{Mobility Edges}} in {{Non-Hermitian Quasicrystals}}}.
\newblock \emph{\bibinfo{journal}{Phys. Rev. Lett.}} \textbf{\bibinfo{volume}{129}}, \bibinfo{pages}{113601} (\bibinfo{year}{2022}).

\bibitem{chuBroadColossalEdge2023}
\bibinfo{author}{Chu, C.-G.} \emph{et~al.}
\newblock \bibinfo{title}{Broad and colossal edge supercurrent in {{Dirac}} semimetal {{Cd3As2 Josephson}} junctions}.
\newblock \emph{\bibinfo{journal}{Nat. Commun.}} \textbf{\bibinfo{volume}{14}}, \bibinfo{pages}{6162} (\bibinfo{year}{2023}).

\bibitem{ochkanNonHermitianTopologyMultiterminal2024}
\bibinfo{author}{Ochkan, K.} \emph{et~al.}
\newblock \bibinfo{title}{Non-{{Hermitian}} topology in a multi-terminal quantum {{Hall}} device}.
\newblock \emph{\bibinfo{journal}{Nat. Phys.}} \textbf{\bibinfo{volume}{20}}, \bibinfo{pages}{395--401} (\bibinfo{year}{2024}).

\bibitem{zhaoTwodimensionalNonHermitianSkin2025}
\bibinfo{author}{Zhao, E.} \emph{et~al.}
\newblock \bibinfo{title}{Two-dimensional non-{{Hermitian}} skin effect in an ultracold {{Fermi}} gas}.
\newblock \emph{\bibinfo{journal}{Nature}} \textbf{\bibinfo{volume}{637}}, \bibinfo{pages}{565--573} (\bibinfo{year}{2025}).

\bibitem{yokomizoNonBlochBandTheory2019}
\bibinfo{author}{Yokomizo, K.} \& \bibinfo{author}{Murakami, S.}
\newblock \bibinfo{title}{Non-{{Bloch Band Theory}} of {{Non-Hermitian Systems}}}.
\newblock \emph{\bibinfo{journal}{Phys. Rev. Lett.}} \textbf{\bibinfo{volume}{123}}, \bibinfo{pages}{066404} (\bibinfo{year}{2019}).

\bibitem{yangNonHermitianBulkBoundaryCorrespondence2020}
\bibinfo{author}{Yang, Z.}, \bibinfo{author}{Zhang, K.}, \bibinfo{author}{Fang, C.} \& \bibinfo{author}{Hu, J.}
\newblock \bibinfo{title}{Non-{{Hermitian Bulk-Boundary Correspondence}} and {{Auxiliary Generalized Brillouin Zone Theory}}}.
\newblock \emph{\bibinfo{journal}{Phys. Rev. Lett.}} \textbf{\bibinfo{volume}{125}}, \bibinfo{pages}{226402} (\bibinfo{year}{2020}).

\bibitem{okumaTopologicalPhaseTransition2019}
\bibinfo{author}{Okuma, N.} \& \bibinfo{author}{Sato, M.}
\newblock \bibinfo{title}{Topological {{Phase Transition Driven}} by {{Infinitesimal Instability}}: {{Majorana Fermions}} in {{Non-Hermitian Spintronics}}}.
\newblock \emph{\bibinfo{journal}{Phys. Rev. Lett.}} \textbf{\bibinfo{volume}{123}}, \bibinfo{pages}{097701} (\bibinfo{year}{2019}).

\bibitem{zhuPhotonicNonHermitianSkin2020}
\bibinfo{author}{Zhu, X.} \emph{et~al.}
\newblock \bibinfo{title}{Photonic non-{{Hermitian}} skin effect and non-{{Bloch}} bulk-boundary correspondence}.
\newblock \emph{\bibinfo{journal}{Phys. Rev. Res.}} \textbf{\bibinfo{volume}{2}}, \bibinfo{pages}{013280} (\bibinfo{year}{2020}).

\bibitem{bartlettIlluminatingBulkboundaryCorrespondence2021}
\bibinfo{author}{Bartlett, J.}, \bibinfo{author}{Hu, H.} \& \bibinfo{author}{Zhao, E.}
\newblock \bibinfo{title}{Illuminating the bulk-boundary correspondence of a non-{{Hermitian}} stub lattice with {{Majorana}} stars}.
\newblock \emph{\bibinfo{journal}{Phys. Rev. B}} \textbf{\bibinfo{volume}{104}}, \bibinfo{pages}{195131} (\bibinfo{year}{2021}).

\bibitem{longhiNonHermitianTopologicalPhase2021}
\bibinfo{author}{Longhi, S.}
\newblock \bibinfo{title}{Non-{{Hermitian}} topological phase transitions in superlattices and the optical {{Dirac}} equation}.
\newblock \emph{\bibinfo{journal}{Opt. Lett.}} \textbf{\bibinfo{volume}{46}}, \bibinfo{pages}{4470--4473} (\bibinfo{year}{2021}).

\bibitem{zhangUniversalNonHermitianSkin2022}
\bibinfo{author}{Zhang, K.}, \bibinfo{author}{Yang, Z.} \& \bibinfo{author}{Fang, C.}
\newblock \bibinfo{title}{Universal non-{{Hermitian}} skin effect in two and higher dimensions}.
\newblock \emph{\bibinfo{journal}{Nat. Commun.}} \textbf{\bibinfo{volume}{13}}, \bibinfo{pages}{2496} (\bibinfo{year}{2022}).

\bibitem{chengTruncationdependent$mathcalPT$Phase2022}
\bibinfo{author}{Cheng, D.} \emph{et~al.}
\newblock \bibinfo{title}{Truncation-dependent \${\textbackslash}mathcal\{\vphantom\}{{PT}}\vphantom\{\}\$ phase transition for the edge states of a two-dimensional non-{{Hermitian}} system}.
\newblock \emph{\bibinfo{journal}{Phys. Rev. B}} \textbf{\bibinfo{volume}{105}}, \bibinfo{pages}{L201105} (\bibinfo{year}{2022}).

\bibitem{fangGeometrydependentSkinEffects2022}
\bibinfo{author}{Fang, Z.}, \bibinfo{author}{Hu, M.}, \bibinfo{author}{Zhou, L.} \& \bibinfo{author}{Ding, K.}
\newblock \bibinfo{title}{Geometry-dependent skin effects in reciprocal photonic crystals}.
\newblock \emph{\bibinfo{journal}{Nanophotonics}} \textbf{\bibinfo{volume}{11}}, \bibinfo{pages}{3447--3456} (\bibinfo{year}{2022}).

\bibitem{wanObservationGeometrydependentSkin2023}
\bibinfo{author}{Wan, T.}, \bibinfo{author}{Zhang, K.}, \bibinfo{author}{Li, J.}, \bibinfo{author}{Yang, Z.} \& \bibinfo{author}{Yang, Z.}
\newblock \bibinfo{title}{Observation of the geometry-dependent skin effect and dynamical degeneracy splitting}.
\newblock \emph{\bibinfo{journal}{Sci. Bull.}} \textbf{\bibinfo{volume}{68}}, \bibinfo{pages}{2330--2335} (\bibinfo{year}{2023}).

\bibitem{zhouObservationGeometrydependentSkin2023}
\bibinfo{author}{Zhou, Q.} \emph{et~al.}
\newblock \bibinfo{title}{Observation of geometry-dependent skin effect in non-{{Hermitian}} phononic crystals with exceptional points}.
\newblock \emph{\bibinfo{journal}{Nat. Commun.}} \textbf{\bibinfo{volume}{14}}, \bibinfo{pages}{4569} (\bibinfo{year}{2023}).

\bibitem{qinGeometrydependentSkinEffect2024}
\bibinfo{author}{Qin, Y.}, \bibinfo{author}{Zhang, K.} \& \bibinfo{author}{Li, L.}
\newblock \bibinfo{title}{Geometry-dependent skin effect and anisotropic {{Bloch}} oscillations in a non-{{Hermitian}} optical lattice}.
\newblock \emph{\bibinfo{journal}{Phys. Rev. A}} \textbf{\bibinfo{volume}{109}}, \bibinfo{pages}{023317} (\bibinfo{year}{2024}).

\bibitem{wangAmoebaFormulationNonBloch2024}
\bibinfo{author}{Wang, H.-Y.}, \bibinfo{author}{Song, F.} \& \bibinfo{author}{Wang, Z.}
\newblock \bibinfo{title}{Amoeba {{Formulation}} of {{Non-Bloch Band Theory}} in {{Arbitrary Dimensions}}}.
\newblock \emph{\bibinfo{journal}{Phys. Rev. X}} \textbf{\bibinfo{volume}{14}}, \bibinfo{pages}{021011} (\bibinfo{year}{2024}).


\bibitem{zhangDynamicalDegeneracySplitting2023}
\bibinfo{author}{Zhang, K.}, \bibinfo{author}{Fang, C.} \& \bibinfo{author}{Yang, Z.}
\newblock \bibinfo{title}{Dynamical {{Degeneracy Splitting}} and {{Directional Invisibility}} in {{Non-Hermitian Systems}}}.
\newblock \emph{\bibinfo{journal}{Phys. Rev. Lett.}} \textbf{\bibinfo{volume}{131}}, \bibinfo{pages}{036402} (\bibinfo{year}{2023}).

\bibitem{zhangAlgebraicNonHermitianSkin2025}
\bibinfo{author}{Zhang, K.}, \bibinfo{author}{Shu, C.} \& \bibinfo{author}{Sun, K.}
\newblock \bibinfo{title}{Algebraic {{Non-Hermitian Skin Effect}} and {{Generalized Fermi Surface Formula}} in {{Arbitrary Dimensions}}}.
\newblock \emph{\bibinfo{journal}{Phys. Rev. X}} \textbf{\bibinfo{volume}{15}}, \bibinfo{pages}{031039} (\bibinfo{year}{2025}).


\bibitem{yaoNonHermitianChernBands2018}
\bibinfo{author}{Yao, S.}, \bibinfo{author}{Song, F.} \& \bibinfo{author}{Wang, Z.}
\newblock \bibinfo{title}{Non-{{Hermitian Chern Bands}}}.
\newblock \emph{\bibinfo{journal}{Phys. Rev. Lett.}} \textbf{\bibinfo{volume}{121}}, \bibinfo{pages}{136802} (\bibinfo{year}{2018}).

\bibitem{yokomizoNonBlochBandsTwodimensional2023}
\bibinfo{author}{Yokomizo, K.} \& \bibinfo{author}{Murakami, S.}
\newblock \bibinfo{title}{Non-{{Bloch}} bands in two-dimensional non-{{Hermitian}} systems}.
\newblock \emph{\bibinfo{journal}{Phys. Rev. B}} \textbf{\bibinfo{volume}{107}}, \bibinfo{pages}{195112} (\bibinfo{year}{2023}).

\bibitem{jiangDimensionalTransmutationNonHermiticity2023}
\bibinfo{author}{Jiang, H.} \& \bibinfo{author}{Lee, C.~H.}
\newblock \bibinfo{title}{Dimensional transmutation from non-{{Hermiticity}}}.
\newblock \emph{\bibinfo{journal}{Phys. Rev. Lett.}} \textbf{\bibinfo{volume}{131}}, \bibinfo{pages}{076401} (\bibinfo{year}{2023}).

\bibitem{zhangEdgeTheoryNonHermitian2024}
\bibinfo{author}{Zhang, K.}, \bibinfo{author}{Yang, Z.} \& \bibinfo{author}{Sun, K.}
\newblock \bibinfo{title}{Edge theory of non-{{Hermitian}} skin modes in higher dimensions}.
\newblock \emph{\bibinfo{journal}{Phys. Rev. B}} \textbf{\bibinfo{volume}{109}}, \bibinfo{pages}{165127} (\bibinfo{year}{2024}).

\bibitem{huTopologicalOriginNonHermitian2024}
\bibinfo{author}{Hu, H.}
\newblock \bibinfo{title}{Topological origin of non-{{Hermitian}} skin effect in higher dimensions and uniform spectra}.
\newblock \emph{\bibinfo{journal}{Sci. Bull.}} \textbf{\bibinfo{volume}{70}}, \bibinfo{pages}{51--57} (\bibinfo{year}{2025}).

\end{thebibliography}

\end{document}
